%%=============================================================================
%% Voorwoord
%%=============================================================================

\chapter*{\IfLanguageName{dutch}{Woord vooraf}{Preface}}%
\label{ch:voorwoord}

%% Het voorwoord is het enige deel van de bachelorproef waar je vanuit je
%% eigen standpunt (``ik-vorm'') mag schrijven. Je kan hier bv. motiveren
%% waarom jij het onderwerp wil bespreken.
%% Vergeet ook niet te bedanken wie je geholpen/gesteund/... heeft

Ik koos dit onderwerp vanuit een interesse in security die al langer aanwezig was, maar die tijdens mijn stage bij Devinity een concrete invulling kreeg. Wat me opviel toen ik de werking van het bedrijf beter leerde kennen, was niet zozeer een gebrek aan tools, maar een gebrek aan samenhang tussen die tools. Logs bleven ongelezen, meldingen kwamen binnen zonder prioriteit, en er was geen gestructureerde manier om op incidenten te reageren. Dat is het soort probleem waarbij je als security specialist meteen denkt: ``dit is op te lossen, waarom is dit er nog niet?'' Die combinatie van herkenbaarheid en haalbaarheid maakte het een logische keuze als onderwerp voor mijn bachelorproef.

Ik wil in de eerste plaats mijn promotor Tom Neels bedanken voor de begeleiding doorheen het proces, de feedback op de tussentijdse versies en de vrijheid om zelf keuzes te maken over de richting van het onderzoek. Daarnaast gaat mijn dank uit naar de collega's bij Devinity die tijd vrijmaakten voor gesprekken, vragen beantwoordden over de bestaande toolset en openstonden voor het uitproberen van nieuwe dingen op de QA-omgeving. Zonder die toegang en dat vertrouwen was het proof-of-concept niet mogelijk geweest.

Ten slotte kijk ik met een positief gevoel terug op dit traject, omdat het mij niet alleen technisch heeft uitgedaagd, maar ook een beter inzicht heeft gegeven in hoe securityprocessen er in de praktijk uitzien.

